\paragraph{Balance firmado y suma con cociente entrante.}
\label{ampl-alfa-z-s-aclaracion}
La división posicional consta de dos etapas consecutivas. Con
\(K_j^{\mathrm{per}}=K_{1+((j-1)\bmod12)}\), que designa la
prolongación periódica denotada también por \(K_j\) en
\eqref{eq:alpha-recurrencia}, el balance anterior al transporte del
cociente es
\[
 Z_j:=P_j+E_j-\Phi_j-K_j^{\mathrm{per}}.
\]
La variable $s_j$ de \eqref{eq:alpha-recurrencia} designa la suma que ya
incluye la contribución entrante, por lo que
\[
 s_j=Z_j+c_{j+1},\qquad
 A_j=Z_j+c_{j+1}-1000c_j.
\]
La configuración incidencial
$H^-_\alpha$ utiliza los signos del balance $Z_j$ de la ventana inicial;
la determinación de las cifras utiliza además los cocientes de enlace
$c_{j+1}$ y $c_j$. Esta distinción conserva conjuntamente el soporte
firmado y la normalización de la expansión.
